\documentclass[10pt,letterpaper]{article}
\usepackage{cornell-notes}

\title{Chapter 13: General Family Algorithms}
\author{}
\date{}

\setDocTitle{Chapter 13: General Family Algorithms}
\setDocSubtitle{Combinatorial Algorithms}
\setDocVersion{v1.0}
\setDocStatus{Study Companion}
\setDocOwner{Combinatorial Algorithms Cornell Notes}
\setDocAuthor{}
\setDocDate{\today}
\setDocSummary{Cornell-format study and review notes for Chapter 13: General Family Algorithms.}

\setCornellCollection{Combinatorial Algorithms Cornell Notes}
\setCornellUnitType{Chapter}
\setCornellUnitNumber{13}
\setCornellUnitTitle{General Family Algorithms}
\setCornellCompanionLabel{Study and review companion}

\hypersetup{pdftitle={Chapter 13: General Family Algorithms},pdfsubject={Cornell notes},pdfauthor={}}

\begin{document}
\maketitle
\begin{tcolorbox}[colback=Paper,colframe=Ink]{\small\bfseries\color{Teal}CORNELL NOTES}\par{\LARGE\bfseries Chapter 13: Sequencing, Ranking, and Selection in General Combinatorial Families (SELECT)}\par\medskip\textbf{Topic:} Recursive-family path coding\hfill\textbf{Date:} \today\par\textbf{Study purpose:} Unify next, rank, unrank, and uniform selection.\end{tcolorbox}
\begin{tcolorbox}[colback=Teal!5,colframe=Teal,title=Essential question]How does a recursively decomposable family become a directed path system supporting four universal operations?\end{tcolorbox}
\begin{tcolorbox}[colback=white,colframe=Mist,title=Learning objectives]\begin{itemize}\item Model objects as complete walks to a terminal vertex.\item Derive the object-count recurrence and edge prefix counts.\item Sequence objects by lexicographic edge codes.\item Rank and unrank paths.\item Sample a path uniformly from subtree sizes.\item Distinguish coded objects from decoded familiar representations.\end{itemize}\textbf{Prerequisites:} directed acyclic graphs, recurrences, lexicographic order, weighted sampling.\end{tcolorbox}
\section*{Cornell notes}
\begin{longtable}{@{}Q!{\color{Mist}\vrule width 1pt}A@{}}\cornellhead
\cue{What is a combinatorial family?}{A locally finite directed graph has a unique minimal terminal vertex. Every edge moves downward in a partial order, and every nonterminal vertex has ranked outgoing edges. A complete directed walk from order vertex $v$ to the terminal vertex is one object of order $v$. Parallel edges represent multiple recursive copies.}
\cue{How is an object coded?}{If a path uses edges $e_1,\ldots,e_s$, record each edge's local rank. The resulting code word determines the path when the start vertex is known. Lexicographic order on code words defines the family order.}
\cue{How are objects counted?}{Let $b(v)$ be the number of complete walks from $v$. Then $b(\tau)=1$ at the terminal vertex and
\[
b(v)=\sum_w g_{vw}b(w),
\]
where $g_{vw}$ is the number of edges from $v$ to $w$. For an outgoing edge $e$, let $f(e)$ be the sum of the subtree counts for all earlier outgoing edges.}
\cue{Task 1: sequence}{Starting at the terminal end of the current path, back up to the rightmost edge that is not the last outgoing edge of its vertex. Increment that edge, then complete the suffix by repeatedly taking edge rank zero. If no edge can advance, the current object is final.}
\cue{Task 2: rank}{For a path $e_1,\ldots,e_s$, add the prefix offsets:
\[
r=\sum_{i=1}^{s} f(e_i).
\]
Each offset counts all objects whose first difference from the given path occurs at that position on an earlier edge. Ranks range from $0$ through $b(v)-1$.}
\cue{Task 3: unrank}{Start with residual rank $r$ at $v$. Choose the greatest outgoing edge whose prefix offset does not exceed the residual, subtract that offset, and continue at the edge's endpoint. On reaching the terminal vertex, the selected path is the unique object of rank $r$.}
\cue{Task 4: random selection}{Either choose a uniform integer in $[0,b(v)-1]$ and unrank it, or at each vertex choose successor $w$ with probability $g_{vw}b(w)/b(v)$ and then choose one of the parallel edges uniformly. Products telescope, giving every complete walk probability $1/b(v)$.}
\cue{What does SELECT generalize?}{The formal routine supports seven recurrence families encoded on lattice-point vertices: fixed-size subsets, set partitions with a fixed number of classes, permutations with a fixed number of cycles, finite-field subspaces, permutations with a fixed number of runs, integer partitions with fixed largest part, and fixed-part compositions. Precomputed counts supply the subtree weights.}
\cue{Encoding versus decoding}{SELECT manipulates vertex and edge arrays describing a path. Many consumers need a conventional set, partition, permutation, or composition. Decoding interprets the edges as construction steps. Encoding is the inverse operation and is separate from universal path traversal.}
\cue{Tradeoffs and cautions}{A generic routine gains reuse but may be slower or less compact than a specialized generator. Counts may overflow; edge order is part of the public ranking contract; parallel edges must remain distinct; and changing decoding conventions must not change the underlying code order accidentally.}
\cue{Retrieval practice}{\begin{enumerate}\item What is an object in the graph model?\item What does $b(v)$ count?\item Why does NEXT reset the suffix to zero edges?\item How is rank accumulated?\item Why do random path probabilities telescope?\item What must decoding provide?\end{enumerate}}
\end{longtable}
\begin{tcolorbox}[colback=Ink!4,colframe=Ink,title=Cornell summary]
A recursively decomposable combinatorial family can be represented as a directed graph whose complete walks are objects. Ranked outgoing edges turn each walk into a lexicographic code. Subtree counts then support four operations with one common structure: advance the rightmost movable edge for sequencing, sum preceding subtree sizes for ranking, descend by residual rank for unranking, and choose successors in proportion to subtree sizes for uniform selection. SELECT specializes these rules to several two-parameter families after their counting tables are computed. The framework separates universal path manipulation from family-specific encoding and decoding. Its principal value is conceptual and reusable correctness; specialized routines may still offer better state locality, arithmetic stability, or transition cost.
\end{tcolorbox}
\end{document}
